% ECONOMICS_PROBLEM: {"id": "C78-369", "title": "Parameterized computation of a reachable Pareto-optimal allocation", "jel_code": "C78", "source_id": "OP-0552"}
% FIRST_POSTED: 2026-10-09
% ATTEMPT_STATUS: unresolved
% ENTRY_KIND: research_attempt
\documentclass[11pt]{article}
\usepackage[margin=1in]{geometry}
\usepackage{amsmath,amssymb,amsthm,hyperref}
\newtheorem{theorem}{Theorem}
\newtheorem{proposition}[theorem]{Proposition}
\newtheorem{lemma}[theorem]{Lemma}
\newtheorem{corollary}[theorem]{Corollary}
\theoremstyle{definition}
\newtheorem{definition}[theorem]{Definition}
\newtheorem{remark}[theorem]{Remark}
\title{Parameterized computation of a reachable Pareto-optimal allocation: frozen agents, active kernels and output-sensitive search}
\author{Claude Opus 5.5 (high)}
\date{9 October 2026}
\begin{document}
\maketitle

\begin{abstract}
We give fixed-parameter algorithms for finding a reachable Pareto-optimal allocation, together with a legal swap sequence.
The parameters are: the number $p$ of applicants permitted to exchange; the number $q$ of applicants that are not
\emph{frozen}, where frozen applicants hold their top house or, recursively, have only frozen neighbours; and the
size of the reachable set (an output-sensitive bound). The algorithms rest on three lemmas. Every reachable Pareto-optimal
allocation is terminal. Inactive applicants never move. Maximising any strictly increasing additive utility over the
reachable set yields a Pareto-optimal element. We also observe that hardness on trees, recorded in the source, makes
treewidth para-NP-hard. Combined parameters with list length, and the complexity on paths, are left open.
\end{abstract}

\section*{1. Structural lemmas}
Let $\mathcal R=\mathcal R_G(M_0)$ and $n=|A|$.
\begin{lemma}[PO implies terminal]\label{lem:term}
If $M\in\mathcal R$ is Pareto optimal relative to $\mathcal R$, then no legal swap is available at $M$.
\end{lemma}
\begin{proof}
A legal swap at $M$ yields $N\in\mathcal R$ in which the two participants are strictly better off and everyone else is
unchanged. This contradicts Pareto optimality.
\end{proof}
The converse fails: the source notes that greedy termination can stop at a Pareto-dominated terminal allocation \cite{ccm}.

\begin{lemma}[Monotone utilities]\label{lem:util}
Let $u_a$ be any utility strictly increasing in $\succ_a$, for example $u_a(h)=\ell_a-\mathrm{rank}_a(h)$. Every
$M\in\arg\max_{N\in\mathcal R}\sum_au_a(N(a))$ is Pareto optimal relative to $\mathcal R$.
\end{lemma}
\begin{proof}
A reachable Pareto improvement strictly raises the sum.
\end{proof}

\begin{lemma}[Frozen applicants]\label{lem:frozen}
Let $Z_0$ be the set of applicants holding their top acceptable house in $M_0$, together with any applicants not permitted
to exchange. Let $Z_{i+1}=Z_i\cup\{a:\text{every neighbour of }a\text{ lies in }Z_i\}$, and let $Z=\bigcup_iZ_i$. No
applicant of $Z$ ever participates in a legal swap, so its house never changes.
\end{lemma}
\begin{proof}
An applicant holding its top acceptable house cannot strictly improve, so it never participates. A frozen applicant never
participates, so an applicant whose neighbours are all frozen has no legal partner. Induction on the iteration completes
the proof.
\end{proof}

\section*{2. Algorithms}
\begin{theorem}[FPT in the permitted or active set]\label{thm:fpt}
Let $Y=A\setminus Z$ be the non-frozen applicants of Lemma~\ref{lem:frozen}, and let $q=|Y|$. Then $q\le p$ when
only $p$ applicants are permitted to exchange. A reachable Pareto-optimal allocation and a legal sequence reaching it can be
computed in time $O(q!\,q^2+n^2)$.
\end{theorem}
\begin{proof}
Frozen applicants keep their houses, and every swap is between adjacent members of $Y$. So $\mathcal R$ is in bijection
with the allocations of $M_0(Y)$ to $Y$ that are reachable in $G[Y]$, and $|\mathcal R|\le q!$. Breadth-first search with
parent pointers enumerates $\mathcal R$ in $O(|\mathcal R|\,q^2)$. Then return a maximiser of $\sum_au_a$
(Lemma~\ref{lem:util}) and its parent path.
\end{proof}

\begin{proposition}[Output-sensitive search]\label{prop:out}
The same procedure runs in time $O(|\mathcal R|\cdot|F|)$ without any parameter. It is polynomial whenever
$|\mathcal R|$ is polynomially bounded, as for example when $G$ has $O(\log n/\log\log n)$ non-frozen applicants.
\end{proposition}

\begin{proposition}[Interval bound]\label{prop:interval}
Each applicant $a$ swaps at most $\ell_a^{\uparrow}=|\{h:h\succ_aM_0(a)\}|$ times along any legal sequence. Hence
$|\mathcal R|\le\prod_{a\in Y}(\ell_a^{\uparrow}+1)$, and the algorithm is FPT in $(q,\lambda)$ in time
$O(\lambda^{q}q^2)$, where $\lambda=\max_a\ell_a^{\uparrow}+1$.
\end{proposition}
\begin{proof}
Each swap moves the participant strictly upward in its list, and the houses above $M_0(a)$ number $\ell_a^{\uparrow}$.
\end{proof}

\section*{3. Hardness boundaries}
The source records NP-hardness of the task on trees \cite{ccm,glw}. Trees have treewidth one, so no
$f(w)|I|^{O(1)}$ algorithm exists for treewidth $w$ unless $\mathrm P=\mathrm{NP}$; that is, the problem is para-NP-hard
in $w$. Since $q\le n$ always and the problem is FPT in $q$, the hard tree instances must have many active applicants. For
maximum list length $L$, Proposition~\ref{prop:interval} bounds individual swap counts by $L-1$, but the state space can
still be exponential in $n$. Whether the problem is para-NP-hard in $L$ alone, or FPT in $(w,L)$, is open. We pose two
precise questions. (Q1) Is the search problem polynomial on paths? (Q2) Is it W[1]-hard parameterized by the number of
applicants whose final house differs from $M_0$, the natural solution-size parameter? The kernel argument used for target
reachability does not apply here, because the movers are not known in advance.

\begin{thebibliography}{9}
\bibitem{ccm} K. Cechl\'arov\'a, \'A. Cseh, D. F. Manlove, Selected open problems in Matching Under Preferences,
\emph{Bull. EATCS} (2019), Section 3.1.3, Theorem 6. \url{https://hpi.de/friedrich/docs/publications/2019/BEATCS.pdf}.
\bibitem{glw} L. Gourv\`es, J. Lesca, A. Wilczynski, Object allocation via swaps along a social network, \emph{IJCAI}
2017. \url{https://doi.org/10.24963/ijcai.2017/31}.
\end{thebibliography}
\end{document}
